\chapter{Transparência Empírica e Tratamento de Variáveis} % Main appendix title
\label{apendicea:metodologia}

Esta seção técnica documenta, de forma exaustiva e rigorosa, o processo de seleção, mensuração e imputação das variáveis utilizadas na modelagem preditiva desta dissertação (focando nas métricas socio-clínicas e operacionais).

Em consonância com as boas práticas de Ciência de Dados em Saúde e visando garantir a reprodutibilidade dos achados preditivos apresentados no Capítulo 5, realizamos um diagnóstico estatístico preliminar de cada característica dos pacientes em relação ao desfecho prognóstico principal: a incapacidade física instalada (\textit{Grau 2 de Incapacidade Física - G2D}) \cite{rocha_transparencia_2023, silva_cura_2024}.

\section{Análise Univariada e Imputação}

O dataset do SINAN frequentemente sofre com dados incompletos (\textit{missing values}). Para mitigar possíveis ruídos algorítmicos decorrentes de atributos enviesados, calculamos a completude histórica percentual.

A estratégia de tratamento obedeceu aos critérios de robustez estatística: variáveis categóricas incompletas foram imputadas utilizando sua \textit{Moda} intra-cluster, enquanto variáveis numéricas com distorção de cauda pesada (altos \textit{outliers}) foram imputadas via \textit{Mediana} estatística \cite{pereira2021imputation}. Os \textit{outliers} foram triados matematicamente utilizando o método do Intervalo Interquartil (IQR).

% Tabela 1: Métricas Descritivas e Tratamento de Variáveis
\begin{table}[H]
\centering
\caption{Métricas Descritivas e Tratamento de Variáveis (SINAN)}
\label{tab:missing_data}
\begin{tabular}{lrlll}
\toprule
Variavel & Missing(\%) & Outliers IQR(\%) & Metodo de Imputacao & Media/Moda \\
\midrule
CLASSOPERA & 0.120000 & 0.160000 & Mediana (Robust in Skewed) [210] & 1.670000 \\
CS\_SEXO & 0.000000 & - & Moda (Categorical) [210] & M \\
CS\_RACA & 7.890000 & 2.860000 & Mediana (Robust in Skewed) [210] & 3.040000 \\
CS\_ESCOL\_N & 6.990000 & - & Moda (Categorical) [210] & 01 \\
NU\_LESOES & 3.780000 & 8.150000 & Mediana (Robust in Skewed) [210] & 5.800000 \\
NERVOSAFET & 46.710000 & 3.820000 & Mediana (Robust in Skewed) [210] & 1.580000 \\
DOSE\_RECEB & 33.620000 & 1.590000 & Mediana (Robust in Skewed) [210] & 9.080000 \\
\bottomrule
\end{tabular}

\end{table}

A tabela acima resume os percentuais de completude e a aplicação dos métodos determinísticos para as principais colunas clínicas de nossa amostra modelada. Note-se a ausência substancial na \textit{Dose Recebida} e alto índice de \textit{Outliers IQR} em contagens quantitativas contínuas (ex: número de lesões nervosas afetadas).

\section{Associação e Relevância Preditiva com G2D}

Após o processo de higienização do esquema de dados, aplicamos técnicas de \textit{Feature Selection} para isolar numericamente os fatores que impulsionariam os modelos de aprendizado de máquina em direção ao alvo "Gravidade Diagnóstica" (G2D=1).

Empregamos o \textit{Kruskal-Wallis H-test} \cite{kruskal1952use}, um método não-paramétrico de análise de variância, para verificar se havia evidência estatística sólida de que os níveis/grupos de cada variável categórica apresentavam distribuições diferentes entre si perante o Grau 2. Em paralelo, capturamos o \textit{Mutual Information (MI)} de cada combinação linear codificada (via \textit{One-Hot Encoding}) para capturar relações complexas (non-lineares).

\begin{table}[H]
\centering
\caption{Análise de Variância (Kruskal-Wallis) para G2D}
\label{tab:kw_g2d}
\begin{tabular}{lrrl}
\toprule
Variavel & KW\_Stat & P\_Value & Associacao\_Significativa \\
\midrule
CLASSOPERA & 130.2427 & 0.0000 & Sim \\
CS\_SEXO & 28.9440 & 0.0000 & Sim \\
CS\_RACA & 84.8525 & 0.0000 & Sim \\
CS\_ESCOL\_N & 297.6595 & 0.0000 & Sim \\
NU\_LESOES & 373.4432 & 0.0000 & Sim \\
NERVOSAFET & 193.1721 & 0.0000 & Sim \\
DOSE\_RECEB & 388.7638 & 0.0000 & Sim \\
\bottomrule
\end{tabular}

\end{table}

\begin{table}[H]
\centering
\caption{Relevância Preditiva via Mutual Information (MI)}
\label{tab:mi_g2d}
\begin{tabular}{lr}
\toprule
Feature Encode & Mutual\_Info[360] \\
\midrule
CLASSOPERA & 0.0051 \\
NU\_LESOES & 0.0044 \\
DOSE\_RECEB & 0.0043 \\
CS\_ESCOL\_N\_5 & 0.0043 \\
CS\_ESCOL\_N\_1 & 0.0042 \\
CS\_ESCOL\_N\_3 & 0.0037 \\
CS\_ESCOL\_N\_02 & 0.0037 \\
NERVOSAFET & 0.0036 \\
CS\_ESCOL\_N\_2 & 0.0034 \\
CS\_SEXO\_M & 0.0032 \\
\bottomrule
\end{tabular}

\end{table}

A Tabela \ref{tab:kw_g2d} exibe a estatística de teste \textit{Kruskal-Wallis} (associada a um P-Value). 
Já a Tabela \ref{tab:mi_g2d} demonstra a contribuição preditiva informacional destas mesmas variáveis quando transformadas via codificação \textit{One-Hot} para o processamento de Máquinas de Vetores de Suporte ou \textit{Gradient Boosting} \cite{scikitlearn_MI_2024}.

\vspace{1cm}
O conjunto dessas duas metodologias garante a ancoragem clínica: atestamos que a construção dos modelos multivariados descritos no Capítulo 5 processa uma série de variáveis cujas relevâncias ao desenvolvimento de G2D fogem estritamente do acaso ($p-value \ll 0.05$).
